\documentclass[a4paper,12pt]{report}
\usepackage[serie,niveau={1MA1},nfiche={}, annee={Emilie-Gourd, 2024--2025}, auteur={ns},theme={Rappel sur les intervalles -- Fonctions}]{packages/bandeaux}
\usepackage{packages/intervalles}
\usepackage{packages/boites}
\usepackage{packages/fontes}
\usepackage{packages/courslatex}
\input{pics/cuboid}
\RequirePackage{icomma}
%chemin vers les sources d'exercices
\renewcommand{\path}{./}

\input{./preambules/print_exo.tex}

% Custom label using TikZ
\newcommand{\tikzlabel}[1]{
  \begin{tikzpicture}[remember picture,overlay]
    \node[draw, circle, line width=0.4pt, inner sep=2pt, anchor=east] {#1};
  \end{tikzpicture}
}


\begin{document}
\textLigne{Intervalles -- Rappels}
% Représenter graphiquement les intervalles suivants:
\insertexo{1M-v2rv8}{false}{exo}{true}

% Donner l'écriture mathématique des intervalles suivants:
\insertexo{1M-6ugwm}{false}{exo}{true}

% Traduire en intervalles les phrase suivantes
\insertexo{1M-61ve6}{false}{exo}{true}

% Décrire par une phrase les intervalles suivants
\insertexo{1M-nw439}{false}{exo}{true}

% Calculer les intersections, unions et différences des intervalles suivants
\insertexo{1M-fq51r}{false}{exo}{true}
\insertexo{1M-u14s6}{false}{exo}{true}


\end{document}


